% sec09_1.tex — Section 9.1: Introduction
\section{Introduction}
\label{sec:greens-intro}

In Chapters~2 and~7 we solved homogeneous partial differential equations with homogeneous boundary conditions by the method of separation of variables.  In Chapter~8 we developed the method of eigenfunction expansion to solve nonhomogeneous problems (both nonhomogeneous PDEs and nonhomogeneous boundary conditions).  In this chapter we study \textbf{Green's functions} for time-independent problems.  Green's functions are the fundamental building blocks for solving nonhomogeneous differential equations, both ordinary and partial.  The Green's function $G(\vec{x},\vec{x}_0)$ represents the response at $\vec{x}$ due to a concentrated source at $\vec{x}_0$.

We begin by motivating the idea of a Green's function from the one-dimensional heat equation.
